By (1.11.3), we can write
@@M0@@
with @@M1@@. We take @@M2@@, which is indeed @@M3@@ if @@M4@@. We then obtain:
@@M5@@

For @@M6@@, @@M7@@, so by (1.11.2):
@@M8@@

Consequently, relation (1.11.4) reduces to:
@@M9@@

Now @@M10@@ has rank @@M11@@, so @@M12@@ is the class of an invertible sheaf, and is therefore invertible in @@M13@@. Consequently, relation (1.11.4) reduces to relation (1.11.2) for @@M14@@.

Thus relation (1.11.2) holds for @@M15@@; for @@M16@@, it gives:
@@M17@@
Since @@M18@@, we obtain @@M19@@. We thus show by descending induction on @@M20@@ that the @@M21@@ are all zero, using relations (1.11.2) for @@M22@@ ranging from @@M23@@ to @@M24@@. This shows that @@M25@@ are independent over @@M26@@, and completes the proof of Theorem 1.1.

\sgaStableTarget{sga6.e06.rem.1-12}\textbf{Remark 1.12.} For every @@M27@@-Module locally free of finite type @@M28@@, let @@M29@@ be the class in @@M30@@ of the @@M31@@-th symmetric power of @@M32@@; put @@M33@@; relations (1.11.3) can also be written:
